\section{结论}
\label{sec:conclusion}

总之，我们讨论了格点上准 LF 关联的重正化与匹配中的若干进一步细节。我们提出了分别处理短距离与长距离关联的混合重正化方案：短距离关联可以通过把同一关联子夹在不同外部态之间作比值来重正化，而长距离关联则用 Wilson 线质量重正化来处理，并以一个连续性条件与短距离区域匹配。这样，我们在重正化阶段避免了在大距离引入额外的非微扰效应。我们还提出如何利用关联的渐近长程行为，向超出格点模拟能力范围的大准 LF 距离作外推，从而避免随后的傅里叶变换中的截断。最后，我们把应用于 LaMET 数据时的大动量展开与 CSF 方法作了比较：前者是提取 PDF $x$ 依赖的系统展开，后者则不是。
我们在此的提议有望大幅改进当前 LaMET 格点应用中的计算策略。
